\documentclass[12pt]{article}
\usepackage[utf8]{inputenc}
\usepackage[T1]{fontenc}
\usepackage[margin=1in]{geometry}
\usepackage{setspace}
\usepackage{xcolor}
\usepackage[colorlinks=true, urlcolor=blue!60!black, linkcolor=blue!60!black]{hyperref}
\setlength{\parindent}{2em}
\setlength{\parskip}{1.25em}
\renewcommand{\baselinestretch}{1.0}
\begin{document}
\begin{center}
{\LARGE \textit{Washington v. Glucksberg}, 521 U.S. 702 (1997)\par}
\vspace{0.8em}
{\large Case Brief --- Four-Page Edition\par}
\end{center}
\vspace{1em}

\section{Memory Jogger}

Because an asserted right to physician-assisted suicide is not ``deeply rooted in this Nation's history and tradition,'' Washington's criminal ban on aiding a suicide does not violate the Due Process Clause (\href{https://supreme.justia.com/cases/federal/us/521/702/}{Justia}).

\section{Detailed Case Facts}

Washington had criminalized assisting suicide since 1854. Under its 1975 code, ``[a] person is guilty of promoting a suicide attempt when he knowingly causes or aids another person to attempt suicide,'' a felony punishable by up to five years in prison. Its Natural Death Act of 1979, by contrast, let adult patients direct the withdrawal of life-sustaining treatment, declaring that such withdrawal ``shall not, for any purpose, constitute a suicide,'' but expressly excluded physician-assisted suicide after voters rejected Initiative 119 in 1991.

In 1994 four physicians, including Dr. Harold Glucksberg, who would have prescribed lethal medication to competent, terminally ill patients but for the statute, joined three gravely ill patients (all of whom died during the litigation) and the nonprofit Compassion in Dying in seeking a declaration that the ban violated the Fourteenth Amendment.

\section{Procedural History}

Relying on \textit{Casey} and \textit{Cruzan}, the district court held that the ban imposed an ``undue burden'' on a protected liberty and also violated equal protection. A Ninth Circuit panel reversed, but the en banc court affirmed, recognizing a liberty interest in ``controlling the time and manner of one's death.'' The Supreme Court heard argument on January 8, 1997, with its companion case \textit{Vacco v. Quill}, and reversed on June 26, 1997.

\section{Judicial Votes}

Reversed and remanded, \textbf{9--0}.
\begin{itemize}
\item \textbf{Majority:} Chief Justice Rehnquist, joined by Justices O'Connor, Scalia, Kennedy, and Thomas.
\item \textbf{Concurring:} Justice O'Connor (substantially endorsed by Justices Ginsburg and Breyer).
\item \textbf{Concurring in the judgment:} Justices Stevens, Souter, Ginsburg, and Breyer. There were no dissents.
\end{itemize}

\section{Holding}

The ban is valid on its face and as applied to competent, terminally ill adults. Substantive due process protects only fundamental rights that are, objectively, ``deeply rooted in this Nation's history and tradition'' and ``implicit in the concept of ordered liberty,'' and requires a ``careful description'' of the asserted interest. Properly framed as ``a right to commit suicide which itself includes a right to assistance in doing so,'' the claimed liberty is not fundamental, and the ban is rationally related to legitimate interests. Footnote 24 left open a more particularized as-applied challenge, and five Justices declined to treat the history-and-tradition test as exclusive.

\section{Analysis of Opinions}

\textbf{Majority (Rehnquist).} A 700-year survey from Bracton and Blackstone through modern statutes showed consistent condemnation of assisting suicide, with no exception for the dying. Rejecting the Ninth Circuit's ``right to die'' framing, the Court distinguished \textit{Cruzan} as resting on the common-law rule that unwanted medical touching is a battery, and read \textit{Casey}'s ``mystery of human life'' passage as no license to protect every intimate decision. The ban rationally served five interests: preserving life, preventing suicide, protecting medical ethics, shielding the poor, elderly, and disabled from coercion, and avoiding a slide toward euthanasia, citing Dutch data.

\textbf{O'Connor (concurring).} She reserved whether a competent person in great suffering has an interest in controlling imminent death; both States permitted palliative care even if it hastened death, and the ``laboratory'' of the States should strike the balance.

\textbf{Stevens.} Invoking Donne's ``No man is an island,'' he accepted the facial ruling but stressed that as-applied claims remain open; a dying patient has an interest in deciding \textit{how}, not \textit{whether}, death occurs.

\textbf{Souter.} Applying Harlan's \textit{Poe v. Ullman} ``reasoned judgment'' approach, he found the interest serious but outweighed for now by risks of abuse, which legislatures are better able to investigate.

\textbf{Ginsburg and Breyer.} Ginsburg concurred for O'Connor's reasons; Breyer suggested a ``right to die with dignity'' and warned that a law blocking adequate palliative care could require revisiting the question.

\section{Examples: Future Applications}

\begin{enumerate}
\item \textbf{Same-Side (No Right): Experimental Drugs.} Carefully described, a right to unapproved investigational drugs has no deep historical roots, and safety interests sustain the restriction.
\item \textbf{Same-Side (No Right): Active Euthanasia.} A physician's lethal injection at a paralyzed patient's request lies even further outside tradition than assisted suicide.
\item \textbf{Opposite-Side (Protected): Forced Ventilation.} Overriding a competent patient's refusal of a ventilator and feeding tube infringes the common-law right recognized in \textit{Cruzan} and affirmed in \textit{Glucksberg}.
\item \textbf{Opposite-Side (Likely Invalid): Banning Palliative Sedation.} Five Justices relied on the availability of pain relief that may hasten death; O'Connor and Breyer signaled that forcing severe pain would raise a grave question.
\item \textbf{Fence-Sitter: Surrogacy and Embryo Bans.} Under ``careful description,'' new reproductive technologies lack historical pedigree; framed broadly as procreative autonomy, they might warrant heightened scrutiny.
\end{enumerate}

\section{Critique}

\textbf{Progressive critiques.} \textit{Dobbs} relied on \textit{Glucksberg}'s test to overrule \textit{Roe} and \textit{Casey}, sharpening the debate. Erwin Chemerinsky called the decision ``tragically wrong'' for ignoring whether the State may force a dying patient to endure added suffering, and critics say the method freezes liberty in past eras' assumptions, in tension with \textit{Lawrence} and \textit{Obergefell}.

\textbf{Restraint perspectives.} Originalists and conservatives praise the test for checking judicial policymaking and leaving medical aid in dying to democratic debate, as in Oregon. Yale Kamisar defended deference because assisted-suicide bans do not target a powerless group; others note the decision preserved \textit{Griswold}, \textit{Loving}, and \textit{Casey}.

\section{Key Quotations}

\begin{quote}
``Our established method of substantive-due-process analysis has two primary features: First, we have regularly observed that the Due Process Clause specially protects those fundamental rights and liberties which are, objectively, `deeply rooted in this Nation's history and tradition,' \ldots\ and `implicit in the concept of ordered liberty,' such that `neither liberty nor justice would exist if they were sacrificed,' \textit{Palko v. Connecticut}, 302 U. S. 319, 325, 326 (1937). Second, we have required in substantive-due-process cases a `careful description' of the asserted fundamental liberty interest.''
\end{quote}
\noindent --- Chief Justice Rehnquist, for the Court, 521 U.S. at 720--721.

\begin{quote}
``That many of the rights and liberties protected by the Due Process Clause sound in personal autonomy does not warrant the sweeping conclusion that any and all important, intimate, and personal decisions are so protected, \textit{San Antonio Independent School Dist. v. Rodriguez}, 411 U. S. 1, 33--35 (1973), and \textit{Casey} did not suggest otherwise.''
\end{quote}
\noindent --- Chief Justice Rehnquist, for the Court, 521 U.S. at 727--728.

\begin{quote}
``Throughout the Nation, Americans are engaged in an earnest and profound debate about the morality, legality, and practicality of physician-assisted suicide. Our holding permits this debate to continue, as it should in a democratic society.''
\end{quote}
\noindent --- Chief Justice Rehnquist, for the Court, 521 U.S. at 735.

\end{document}
